%%%%%%%%%%%%%%%%%%%%%%%%%%%%%%%%%%%%%%%%%%%%%%%%%%%%%%%%%%%%%%%%%%%%%%%%%%%
% Chapter: Voltage and Current Sources
% Falstad examples: examples/sources (generated by falstad/circuits/sources.py)
%%%%%%%%%%%%%%%%%%%%%%%%%%%%%%%%%%%%%%%%%%%%%%%%%%%%%%%%%%%%%%%%%%%%%%%%%%%

%%%%%%%%%%%%%%%%%%%%%%%%%%%%%%%%%%%%%%%%%%%%%%%%%%%%%%%%%%%%%%%%%%%%%%%%%%%
\section{Ideal and Real Voltage Sources}
%%%%%%%%%%%%%%%%%%%%%%%%%%%%%%%%%%%%%%%%%%%%%%%%%%%%%%%%%%%%%%%%%%%%%%%%%%%
\nsl{Every circuit needs energy, and every measurement needs a quantity
that is known exactly: a supply voltage, a reference voltage or an
excitation current for a sensor. This chapter shows how real sources
differ from the ideal ones of circuit theory, and how simple circuits with
Zener diodes, transistors and op-amps produce stable voltages and
currents. In the capstone task the bridge needs a stable supply, the
reference ladder of the flash converter needs a stable \SI{4}{\volt}, and
the resistive sensor must not carry more than \SI{1}{\milli\ampere}. The
chapter on voltage references continues with precise references (bandgap
references); here we treat the basic circuits.\newpage}

\begin{description}
\item[{\rot\bf Ideal and real voltage source}]~\\[-\lblskip]
\begin{center}
\begin{circuitikz}[scale=1, transform shape]
  \draw (0,0) to[V, invert, l=$V_0$] (0,3) to[short, -o] (1.5,3);
  \draw (0,0) to[short, -o] (1.5,0);
  \node[font=\small] at (0.6,-0.6) {ideal};
  \draw[dashed, rwucyandark] (2.8,-0.3) rectangle (5.6,3.6);
  \draw (4,0) to[V, invert, l=$V_0$] (4,1.8) to[R, l=$R_i$] (4,3.2) -- (5,3.2) to[short, -o] (6.5,3.2);
  \draw (4,0) to[short, -o] (6.5,0);
  \draw (6.5,3.2) -- (8,3.2) to[R, l=$R_L$, i>^=$I_L$] (8,0) -- (6.5,0);
  \draw (6.5,3.2) to[open, v^=$V_L$] (6.5,0);
  \node[font=\small] at (4.6,-0.6) {real source};
\end{circuitikz}
\end{center}
\bi
	\ite ideal: the voltage is $V_0$ for every current
	\ite real: the terminal voltage drops with the load current
\ei
\keyeq{Terminal voltage of a real source}{$\displaystyle V_L = V_0 - R_i\,I_L = V_0\,\frac{R_L}{R_i+R_L}$}
\os{\newpage}

\item[{\rot\bf Examples of internal resistances}]~\\[-\lblskip]
\bi
	\ite AA alkaline cell: $V_0 = \SI{1.5}{\volt}$, $R_i \approx \SIrange{0.1}{0.3}{\ohm}$ (rises as the cell discharges)
	\ite \SI{9}{\volt} block battery: $R_i \approx \SIrange{1}{2}{\ohm}$
	\ite laboratory power supply in voltage mode: some \si{\milli\ohm}, plus the leads
	\ite voltage divider: $R_i = R_1\parallelto R_2$ (chapter Basics)
	\ite logic output: \SIrange{20}{50}{\ohm}
	\ite op-amp output with feedback: well below \SI{1}{\ohm} at low frequencies
\ei
\nsl{The internal resistance is not a resistor one could find inside the
source; it is a model for everything that makes the voltage drop under
load: the chemistry of a battery, the control loop of a power supply, the
output transistors of a logic gate. A real source is therefore described
by two numbers, $V_0$ and $R_i$, which can be measured: the open-circuit
voltage directly, the internal resistance from a second measurement with a
known load, $R_i = (V_0 - V_L)/I_L$. A short circuit to measure $I_\mathrm{sc}$
directly is not a good idea for batteries and power supplies.}
\os{\newpage}

\item[{\rot\bf Load line and operating point}]~\\[-\lblskip]
\begin{center}
\begin{tikzpicture}
\begin{axis}[width=11cm, height=5.6cm, xlabel={$I_L$ / A}, ylabel={$V_L$ / V},
  xmin=0, xmax=4, ymin=0, ymax=5, grid=major, font=\small]
  \addplot[rwuviolet, very thick, domain=0:3.6] {4.5 - 1.25*x};
  \addplot[black, thick, dashed, domain=0:4] {4.5};
  \addplot[rwucyandark, thick, domain=0:0.5] {10*x};
  \addplot[only marks, mark=*, mark size=2pt] coordinates {(0.4,4.0) (1.8,2.25)};
  \addplot[rwucyandark, thick, domain=0:2.4] {1.25*x};
  \node[anchor=west, font=\scriptsize] at (axis cs:0.45,4.15) {$R_L=\SI{10}{\ohm}$: \SI{4.0}{\volt}, \SI{0.4}{\ampere}};
  \node[anchor=west, font=\scriptsize] at (axis cs:1.9,2.45) {$R_L=R_i$: maximum power};
  \node[anchor=south, font=\scriptsize] at (axis cs:3.4,0.1) {$I_\mathrm{sc}=\SI{3.6}{\ampere}$};
  \node[anchor=south, font=\scriptsize] at (axis cs:3.0,4.5) {ideal source};
\end{axis}
\end{tikzpicture}
\end{center}
\bi
	\ite example: $V_0=\SI{4.5}{\volt}$, $R_i=\SI{1.25}{\ohm}$; load lines of $R_L=\SI{10}{\ohm}$ and $R_L=R_i$
	\ite the operating point is the intersection of source line and load line
\ei
\falstad{real_source}
\os{\newpage}

\item[{\rot\bf Power matching}]~\\[-\lblskip]
\begin{center}
\begin{tikzpicture}
\begin{axis}[width=9.5cm, height=5cm, xlabel={$R_L/R_i$}, ylabel={$P_L/P_\mathrm{max}$, $\eta$},
  xmin=0, xmax=5, ymin=0, ymax=1.1, grid=major, font=\small, legend pos=south east,
  legend style={font=\scriptsize}]
  \addplot[rwuviolet, very thick, domain=0:5, samples=80] {4*x/(1+x)^2};
  \addplot[rwucyandark, thick, dashed, domain=0:5, samples=80] {x/(1+x)};
  \legend{$P_L/P_\mathrm{max}$, efficiency $\eta$}
\end{axis}
\end{tikzpicture}
\end{center}
\bi
	\ite load power $P_L = V_0^2\,R_L/(R_i+R_L)^2$ is maximum for $R_L = R_i$: $P_\mathrm{max} = V_0^2/(4R_i)$
	\ite then half of the power is lost in the source: efficiency $\eta = \SI{50}{\percent}$
	\ite power supplies and signal sources for measurement: $R_L \gg R_i$ (voltage matching)
\ei
\nsl{Power matching is important in radio frequency circuits and for
antennas, where the cable impedance must be matched, and for weak energy
sources such as thermoelectric generators. For supplies and in measurement
circuits it is the wrong goal: we want the voltage to stay constant, so
the load resistance must be large compared with the internal resistance.
The same rule appeared for the loaded voltage divider.}
\os{\newpage}
\end{description}

\exercise{Measuring a real voltage source}
A battery pack shows \SI{4.5}{\volt} without load and \SI{4.0}{\volt} with a load of \SI{10}{\ohm}.
\begin{talist}
\ta Calculate the load current and the internal resistance.
\ta Calculate the short-circuit current.
\ta Which load takes the maximum power from the battery? How large is this power and the efficiency?
\ta Check (a) in Falstad and use the slider to find the maximum load power.
\end{talist}

\begin{solution}
\begin{talist}
\ta $I_L = \SI{4.0}{\volt}/\SI{10}{\ohm} = \SI{0.4}{\ampere}$, $R_i = (\SI{4.5}{\volt}-\SI{4.0}{\volt})/\SI{0.4}{\ampere} = \SI{1.25}{\ohm}$.
\ta $I_\mathrm{sc} = V_0/R_i = \SI{4.5}{\volt}/\SI{1.25}{\ohm} = \SI{3.6}{\ampere}$.
\ta $R_L = R_i = \SI{1.25}{\ohm}$: $P_\mathrm{max} = (\SI{4.5}{\volt})^2/(4\cdot\SI{1.25}{\ohm}) = \SI{4.05}{\watt}$,
    $V_L = \SI{2.25}{\volt}$, $\eta = \SI{50}{\percent}$; the battery heats up with the same \SI{4.05}{\watt}.
\ta \falstad{real_source}: \texttt{V\_ideal} $=\SI{4.50}{\volt}$, \texttt{V\_term} $=\SI{4.00}{\volt}$, $I_L = \SI{0.400}{\ampere}$
    (hover over $R_L$). With the slider at $R_L = \SI{1.25}{\ohm}$ the resistor shows \SI{4.05}{\watt}.
\end{talist}
\end{solution}

%%%%%%%%%%%%%%%%%%%%%%%%%%%%%%%%%%%%%%%%%%%%%%%%%%%%%%%%%%%%%%%%%%%%%%%%%%%
\section{Current Sources and Source Transformation}
%%%%%%%%%%%%%%%%%%%%%%%%%%%%%%%%%%%%%%%%%%%%%%%%%%%%%%%%%%%%%%%%%%%%%%%%%%%
\nsl{A current source delivers a given current whatever voltage the load
needs. In the lab it is less familiar than a voltage source, but it is
everywhere inside integrated circuits: it biases transistor stages, it
charges capacitors linearly, and it excites resistive sensors, because the
voltage across the sensor is then directly proportional to its resistance.
\newpage}

\begin{description}
\item[{\rot\bf Ideal and real current source}]~\\[-\lblskip]
\begin{center}
\begin{circuitikz}[scale=1, transform shape]
  \draw (0,0) to[I, l=$I_0$] (0,3) to[short, -o] (1.5,3);
  \draw (0,0) to[short, -o] (1.5,0);
  \node[font=\small] at (0.6,-0.6) {ideal};
  \draw[dashed, rwucyandark] (2.9,-0.3) rectangle (6,3.5);
  \draw (4,0) to[I, l=$I_0$] (4,3) -- (5.3,3) to[short, -o] (6.8,3);
  \draw (5.3,0) to[R, l_=$R_i$] (5.3,3);
  \draw (4,0) to[short, -o] (6.8,0);
  \draw (6.8,3) -- (8.2,3) to[R, l=$R_L$, i>^=$I_L$] (8.2,0) -- (6.8,0);
  \node[font=\small] at (4.7,-0.6) {real source};
\end{circuitikz}
\end{center}
\keyeq{Load current of a real current source}{$\displaystyle I_L = I_0 - \frac{V_L}{R_i} = I_0\,\frac{R_i}{R_i+R_L}$}
\bi
	\ite a good current source has a \textbf{large} internal resistance ($R_i \gg R_L$)
	\ite an ideal current source must never be left open: its voltage would rise without limit
	\ite real current sources work only in a limited voltage range: the \textbf{compliance range}
\ei
\os{\newpage}

\item[{\rot\bf Source transformation}]~\\[-\lblskip]
\begin{center}
\begin{circuitikz}[scale=1, transform shape]
  \draw (0,0) to[V, invert, l=$V_0$] (0,3) to[R, l=$R_i$, -o] (2.6,3);
  \draw (0,0) to[short, -o] (2.6,0);
  \node at (3.6,1.5) {$\Leftrightarrow$};
  \draw (5.2,0) to[I, l=$I_0$] (5.2,3) to[short, -o] (8,3);
  \draw (6.6,0) to[R, l_=$R_i$] (6.6,3);
  \draw (5.2,0) to[short, -o] (8,0);
\end{circuitikz}
\end{center}
\bi
	\ite both have the same open-circuit voltage $V_0$ and the same short-circuit current $I_0 = V_0/R_i$
	\ite therefore they behave identically for \textbf{every} load (Thevenin $\leftrightarrow$ Norton)
	\ite inside, they are different: the power lost in $R_i$ is not the same
	\ite example: \SI{10}{\volt} with \SI{10}{\kilo\ohm} $\Leftrightarrow$ \SI{1}{\milli\ampere} with \SI{10}{\kilo\ohm}; with $R_L=\SI{4.7}{\kilo\ohm}$ both give \SI{3.20}{\volt}
\ei
\falstad{source_transform}
\nsl{Whether a source is called a voltage source or a current source is a
question of the load. A source with $R_i = \SI{10}{\kilo\ohm}$ is a fairly
good voltage source for a \SI{1}{\mega\ohm} load and a fairly good current
source for a \SI{100}{\ohm} load. Source transformation is also a
practical tool for calculations: series connections are easy with voltage
sources, parallel connections with current sources.}
\os{\newpage}
\end{description}

\exercise{Source transformation}
A source consists of \SI{10}{\volt} in series with \SI{10}{\kilo\ohm}. It feeds $R_L = \SI{4.7}{\kilo\ohm}$.
\begin{talist}
\ta Transform it into a current source with a parallel resistor.
\ta Calculate the load voltage with both models.
\ta Calculate the power dissipated in the internal resistance for both models. What do you conclude?
\ta Check (b) in Falstad.
\end{talist}

\begin{solution}
\begin{talist}
\ta $I_0 = \SI{10}{\volt}/\SI{10}{\kilo\ohm} = \SI{1}{\milli\ampere}$ parallel to \SI{10}{\kilo\ohm}.
\ta Voltage source: $V_L = \SI{10}{\volt}\cdot 4.7/14.7 = \SI{3.20}{\volt}$.
    Current source: $V_L = \SI{1}{\milli\ampere}\cdot(\SI{10}{\kilo\ohm}\parallelto\SI{4.7}{\kilo\ohm}) = \SI{1}{\milli\ampere}\cdot\SI{3.20}{\kilo\ohm} = \SI{3.20}{\volt}$.
\ta Voltage source: $I = \SI{0.680}{\milli\ampere}$ flows through $R_i$, $P = \SI{4.63}{\milli\watt}$.
    Current source: $R_i$ sees \SI{3.20}{\volt}, $P = \SI{1.02}{\milli\watt}$.
    The models are equivalent only at the terminals, not inside.
\ta \falstad{source_transform}: \texttt{VL\_thevenin} $=$ \texttt{VL\_norton} $=\SI{3.20}{\volt}$.
\end{talist}
\end{solution}

%%%%%%%%%%%%%%%%%%%%%%%%%%%%%%%%%%%%%%%%%%%%%%%%%%%%%%%%%%%%%%%%%%%%%%%%%%%
\section{The Zener Diode Shunt Regulator}
%%%%%%%%%%%%%%%%%%%%%%%%%%%%%%%%%%%%%%%%%%%%%%%%%%%%%%%%%%%%%%%%%%%%%%%%%%%
\nsl{A Zener diode is a diode that is operated in reverse breakdown. Above
the breakdown voltage the reverse current rises steeply, so the voltage
stays almost constant over a wide current range. Together with a series
resistor this gives the simplest voltage regulator. It is good enough for
small currents and as a reference for the transistor regulator of the next
section.\newpage}

\begin{description}
\item[{\rot\bf Zener diode characteristic}]~\\[-\lblskip]
\begin{center}
\begin{tikzpicture}
\begin{axis}[width=11cm, height=5.6cm, xlabel={$V$ / V}, ylabel={$I$ / mA},
  xmin=-7, xmax=1.5, ymin=-40, ymax=30, grid=major, font=\small]
  \addplot[rwuviolet, very thick, domain=0:0.85, samples=60] {min(30, 5*exp((x-0.75)/0.04))};
  \addplot[rwuviolet, very thick, domain=-5.75:0, samples=80] {max(-40, -5*exp((-x-5.6)/0.035))};
  \addplot[only marks, mark=*, mark size=1.5pt] coordinates {(-5.6,-5)};
  \node[anchor=east, font=\scriptsize] at (axis cs:-5.75,-5) {$V_Z$ at $I_Z=\SI{5}{\milli\ampere}$};
  \node[anchor=west, font=\scriptsize] at (axis cs:-5.55,-30) {breakdown: $V \approx -V_Z$};
  \node[anchor=west, font=\scriptsize] at (axis cs:0.2,22) {forward: \SI{0.7}{\volt}};
\end{axis}
\end{tikzpicture}
\end{center}
\bi
	\ite breakdown voltages \SIrange{2.4}{75}{\volt} in E24 steps (BZX55, BZX79: \SI{0.5}{\watt}; 1N47xx: \SI{1.3}{\watt})
	\ite datasheet: $V_Z$ at a test current (often \SI{5}{\milli\ampere}), dynamic resistance $r_Z = \Delta V/\Delta I$
\ei
\nsl{Below about \SI{5}{\volt} the breakdown is caused by the Zener effect
(tunnelling), above by avalanche multiplication. Diodes around
\SIrange{5}{6}{\volt} have the smallest temperature coefficient and the
smallest dynamic resistance, which is why \SI{5.6}{\volt} and \SI{6.2}{\volt}
are popular reference values. Low voltage Zener diodes have a soft knee:
a \SI{2.7}{\volt} type has a dynamic resistance of some \SI{100}{\ohm}.
The Falstad model is calibrated so that the current is exactly
\SI{5}{\milli\ampere} at the nominal $V_Z$, and it is much stiffer than a
real diode (a few ohms instead of typically \SIrange{10}{40}{\ohm}).}
\os{\newpage}

\item[{\rot\bf Shunt regulator}]~\\[-\lblskip]
\begin{center}
\begin{circuitikz}[scale=1, transform shape]
  \draw (0,0) to[V, invert, l=$V_\mathrm{in}$] (0,3) to[short, i=$I_R$] (1.2,3)
        to[R, l=$R$] (3.4,3) -- (5.6,3);
  \draw (3.4,3) to[short, i_=$I_Z$] (3.4,2.4) to[zDo, l_=$V_Z$] (3.4,0);
  \draw (5.6,3) to[short, i=$I_L$] (5.6,2.4) to[R, l=$R_L$] (5.6,0);
  \draw (0,0) -- (5.6,0);
  \draw (3.4,0) node[ground]{};
  \draw (5.6,3) node[circ]{} node[above]{$\Vout$};
\end{circuitikz}
\end{center}
\bi
	\ite the resistor takes the voltage difference: $I_R = (V_\mathrm{in}-V_Z)/R$
	\ite the Zener takes what the load does not need: $I_Z = I_R - I_L$
	\ite the Zener must keep a minimum current $I_{Z,\min}$ (about \SI{5}{\milli\ampere}, or \SI{10}{\percent} of $I_{Z,\max}$)
\ei
\os{\newpage}

\item[{\rot\bf Designing a shunt regulator}]~\\
\designcard{Zener shunt regulator}{
\begin{enumerate}\itemsep 0pt
	\item worst case for regulation: lowest $V_\mathrm{in}$, largest load current
	$\Rightarrow\ R \le \dfrac{V_{\mathrm{in},\min}-V_Z}{I_{L,\max}+I_{Z,\min}}$, choose the next \textbf{smaller} E-value
	\item worst case for the Zener: highest $V_\mathrm{in}$, no load
	$\Rightarrow\ I_{Z,\max} = \dfrac{V_{\mathrm{in},\max}-V_Z}{R}$, $P_Z = V_Z\,I_{Z,\max}$
	\item power of the resistor: $P_R = (V_{\mathrm{in},\max}-V_Z)^2/R$
\end{enumerate}}
\bi
	\ite example: \SI{12}{\volt}, \SI{5.6}{\volt}, $I_L = \SIrange{0}{20}{\milli\ampere}$:
	$R \le \SI{6.4}{\volt}/\SI{25}{\milli\ampere} = \SI{256}{\ohm}$ $\Rightarrow$ \SI{220}{\ohm}
	\ite the efficiency is poor: the full current flows all the time, even without load
\ei
\falstad{zener_shunt}
\os{\newpage}

\item[{\rot\bf Regulation and ripple}]~\\[-\lblskip]
\begin{center}
\begin{circuitikz}[scale=1, transform shape]
  \draw (0,0) to[V, invert, l=$\Delta V_\mathrm{in}$] (0,2.5) to[R, l=$R$] (3,2.5) -- (5,2.5);
  \draw (3,2.5) to[R, l_=$r_Z$] (3,0);
  \draw (5,2.5) to[R, l=$R_L$] (5,0);
  \draw (0,0) -- (5,0);
  \draw (5,2.5) node[circ]{} node[above]{$\Delta\Vout$};
\end{circuitikz}
\end{center}
\bi
	\ite for small changes the Zener is a resistor $r_Z$: small-signal equivalent circuit
	\ite line regulation: $\displaystyle\frac{\Delta\Vout}{\Delta V_\mathrm{in}} = \frac{r_Z\parallelto R_L}{R + r_Z\parallelto R_L}\approx\frac{r_Z}{R}$
	\ite load regulation: output resistance $R_\mathrm{out} = R\parallelto r_Z\approx r_Z$
	\ite example $r_Z = \SI{20}{\ohm}$, $R=\SI{220}{\ohm}$: \SI{1}{\volt} ripple at the input $\to$ \SI{83}{\milli\volt} at the output
\ei
\nsl{The Zener regulator suppresses input changes only by the ratio
$r_Z/R$, typically a factor of 10. A larger resistor improves this, but
requires a higher input voltage. Better regulation is obtained when the
Zener diode is fed from a current source, or when a feedback loop with an
amplifier is used (series regulator, voltage reference). The output
resistance $r_Z$ means that the output voltage changes with the load
current: for \SI{20}{\milli\ampere} and $r_Z=\SI{20}{\ohm}$ by
\SI{0.4}{\volt}.}
\os{\newpage}
\end{description}

\exercise{Design of a Zener shunt regulator}
A load needs \SI{5.6}{\volt} and draws \SIrange{0}{20}{\milli\ampere}. The input voltage is \SI{12}{\volt}.
Use a \SI{5.6}{\volt} Zener diode with $I_{Z,\min}=\SI{5}{\milli\ampere}$.
\begin{talist}
\ta Calculate the series resistor and choose an E12 value.
\ta Calculate the maximum Zener current and the power of Zener diode and resistor. Which ratings do you choose?
\ta Does the circuit still regulate at full load when the input drops to \SI{11}{\volt}?
\ta Read the output voltage without load and with $R_L=\SI{270}{\ohm}$ in Falstad and calculate $I_Z$.
\end{talist}

\begin{solution}
\begin{talist}
\ta $R \le (\SI{12}{\volt}-\SI{5.6}{\volt})/(\SI{20}{\milli\ampere}+\SI{5}{\milli\ampere}) = \SI{256}{\ohm}$ $\Rightarrow$ \SI{220}{\ohm}.
\ta Without load: $I_{Z,\max} = \SI{6.4}{\volt}/\SI{220}{\ohm} = \SI{29.1}{\milli\ampere}$,
    $P_Z = \SI{5.6}{\volt}\cdot\SI{29.1}{\milli\ampere} = \SI{163}{\milli\watt}$ $\Rightarrow$ \SI{0.5}{\watt} Zener (BZX79C5V6).
    $P_R = (\SI{6.4}{\volt})^2/\SI{220}{\ohm} = \SI{186}{\milli\watt}$ $\Rightarrow$ \SI{0.5}{\watt} resistor.
\ta $I_R = (\SI{11}{\volt}-\SI{5.6}{\volt})/\SI{220}{\ohm} = \SI{24.5}{\milli\ampere}$, $I_Z = \SI{4.5}{\milli\ampere}$:
    just below $I_{Z,\min}$; the voltage drops slightly, but the circuit still works. For a safe margin choose \SI{200}{\ohm} (E24).
\ta \falstad{zener_shunt}\newline \texttt{Vout\_open} $=\SI{5.65}{\volt}$ ($I_Z = I_R = \SI{28.9}{\milli\ampere}$),
    \texttt{Vout\_load} $=\SI{5.61}{\volt}$,\newline $I_L = \SI{5.61}{\volt}/\SI{270}{\ohm} = \SI{20.8}{\milli\ampere}$,
    $I_R = (12-5.61)\,\si{\volt}/\SI{220}{\ohm} = \SI{29.0}{\milli\ampere}$, $I_Z = \SI{8.3}{\milli\ampere}$.
\end{talist}
\end{solution}

%%%%%%%%%%%%%%%%%%%%%%%%%%%%%%%%%%%%%%%%%%%%%%%%%%%%%%%%%%%%%%%%%%%%%%%%%%%
\section{Series Regulators}
%%%%%%%%%%%%%%%%%%%%%%%%%%%%%%%%%%%%%%%%%%%%%%%%%%%%%%%%%%%%%%%%%%%%%%%%%%%
\nsl{The shunt regulator wastes current and cannot deliver large load
currents. A series regulator puts a transistor between input and output.
The transistor carries only the load current, and the Zener diode only has
to deliver the small base current. Integrated linear regulators add an
error amplifier, a bandgap reference and protection circuits; they are the
standard solution for supplies up to about \SI{1}{\ampere}.\newpage}

\begin{description}
\item[{\rot\bf Emitter follower with Zener diode}]~\\[-\lblskip]
\begin{center}
\begin{circuitikz}[scale=1, transform shape]
  \draw (0,0) to[V, invert, l=$V_\mathrm{in}$] (0,4) -- (5,4);
  \draw (2.5,4) to[R, l_=$R_B$] (2.5,2);
  \draw (2.5,2) to[zDo, l_=$V_Z$] (2.5,0);
  \draw (5,2) node[npn](Q){};
  \draw (Q.C) -- (5,4);
  \draw (Q.B) -- (2.5,2);
  \draw (Q.E) -- (5,0.9) -- (6.5,0.9) to[R, l=$R_L$] (6.5,-0.4) -- (6.5,-0.6);
  \draw (0,0) -- (0,-0.6) -- (6.5,-0.6);
  \draw (2.5,0) -- (2.5,-0.6);
  \draw (2.5,-0.6) node[ground]{};
  \draw (6.5,0.9) node[circ]{} node[above]{$\Vout$};
\end{circuitikz}
\end{center}
\keyeq{Emitter follower regulator}{$\displaystyle \Vout = V_Z - V_{BE} \approx V_Z - \SI{0.65}{\volt}\qquad I_B = \frac{I_L}{\beta+1}$}
\bi
	\ite the Zener sees only $I_B$: load changes hardly affect $V_Z$
\ei
\falstad{series_regulator}
\os{\newpage}

\item[{\rot\bf Properties of the emitter follower regulator}]~\\[-\lblskip]
\bi
	\ite output resistance $\approx r_e + (R_B\parallelto r_Z)/(\beta+1)$ with $r_e = V_T/I_E$ (\SI{0.25}{\ohm} at \SI{100}{\milli\ampere}): much smaller than with the Zener alone
	\ite $V_{BE}$ depends on current and temperature ($-\SI{2}{\milli\volt\per\kelvin}$): $\Vout$ is only accurate to about \SI{0.1}{\volt}
	\ite transistor power: $P_Q = (V_\mathrm{in}-\Vout)\cdot I_L$ (heat sink!)
	\ite design of $R_B$: $\dfrac{V_{\mathrm{in},\min}-V_Z}{R_B} \ge I_{B,\max} + I_{Z,\min}$
	\ite no short-circuit protection: a short at the output destroys the transistor
\ei
\nsl{The Falstad example shows the effect: with \SI{470}{\ohm} (about
\SI{10}{\milli\ampere}) the output is \SI{4.95}{\volt}, with \SI{47}{\ohm}
(\SI{104}{\milli\ampere}) it is \SI{4.89}{\volt}. The Zener voltage itself
changes by only \SI{4}{\milli\volt}; most of the change comes from
$V_{BE}$, which grows by $V_T\ln 10 = \SI{60}{\milli\volt}$ for ten times
the current. An op-amp that compares the output with the reference and
drives the transistor removes this error: that is the principle of the
integrated regulator.}
\os{\newpage}

\item[{\rot\bf Integrated fixed regulator 78xx}]~\\[-\lblskip]
\begin{center}
\begin{circuitikz}[scale=1, transform shape]
  \node[draw, thick, text width=2.4cm, align=center, minimum height=1.6cm] (U) at (5,2) {7805};
  \node[font=\scriptsize, anchor=west] at (U.west) {IN};
  \node[font=\scriptsize, anchor=east] at (U.east) {OUT};
  \node[font=\scriptsize, anchor=south] at (U.south) {GND};
  \draw (0,0) to[V, invert, l=$V_\mathrm{in}$] (0,2) -- (U.west);
  \draw (1.6,2) to[C, l=\SI{330}{\nano\farad}] (1.6,0);
  \draw (U.east) -- (9.8,2) to[R, l=$R_L$] (9.8,0);
  \draw (8.2,2) to[C, l_=\SI{100}{\nano\farad}] (8.2,0);
  \draw (U.south) -- (5,0);
  \draw (0,0) -- (9.8,0);
  \draw (5,0) node[ground]{};
  \draw (9.8,2) node[circ]{} node[above]{\SI{5}{\volt}};
\end{circuitikz}
\end{center}
\bi
	\ite 7805, 7812, \ldots: fixed output voltage, \SI{1}{\ampere} (TO-220), \SI{0.1}{\ampere} (78L05); 79xx for negative voltages
	\ite dropout about \SI{2}{\volt}: $V_\mathrm{in} \ge \Vout + \SI{2}{\volt}$; LDO regulators: \SIrange{0.1}{0.5}{\volt}
	\ite internal current limit and thermal shutdown; capacitors close to the pins for stability
\ei
\os{\newpage}

\item[{\rot\bf Adjustable regulator LM317}]~\\[-\lblskip]
\begin{center}
\begin{circuitikz}[scale=1, transform shape]
  \node[draw, thick, text width=2.4cm, align=center, minimum height=1.6cm] (U) at (3.5,3) {LM317};
  \node[font=\scriptsize, anchor=west] at (U.west) {IN};
  \node[font=\scriptsize, anchor=east] at (U.east) {OUT};
  \node[font=\scriptsize, anchor=south] at (U.south) {ADJ};
  \draw (0,0) to[V, invert, l=$V_\mathrm{in}$] (0,3) -- (U.west);
  \draw (U.east) -- (6.5,3) to[R, l=$R_1$] (6.5,1.4) to[R, l=$R_2$] (6.5,0);
  \draw (U.south) -- (3.5,1.4) -- (6.5,1.4);
  \draw (6.5,3) -- (8.8,3) to[R, l=$R_L$] (8.8,0);
  \draw (0,0) -- (8.8,0);
  \draw (6.5,0) node[ground]{};
  \draw (8.8,3) node[circ]{} node[above]{$\Vout$};
\end{circuitikz}
\end{center}
\keyeq{LM317}{$\displaystyle \Vout = \SI{1.25}{\volt}\left(1+\frac{R_2}{R_1}\right) + I_\mathrm{ADJ}\,R_2\qquad I_\mathrm{ADJ}\approx\SI{50}{\micro\ampere}$}
\bi
	\ite the regulator keeps \SI{1.25}{\volt} between OUT and ADJ; $R_1 = \SI{240}{\ohm}$ sets \SI{5.2}{\milli\ampere}, enough for the minimum load
\ei
\nsl{The LM317 has no ground pin. Its internal reference sits between OUT
and ADJ, and all its quiescent current flows out of the output. The
current through $R_1$ ($\SI{1.25}{\volt}/R_1$) also flows through $R_2$;
together with the small $I_\mathrm{ADJ}$ this gives the formula. The data
sheet requires a minimum load current of about \SI{10}{\milli\ampere}
(typically \SI{3.5}{\milli\ampere}); the divider takes care of it. Without
$R_2$ and with the load in place of $R_2$, the LM317 becomes a current
source: $I = \SI{1.25}{\volt}/R_1$. Falstad has no model for the LM317;
its principle (reference, error amplifier, series transistor) is the topic
of the chapter on voltage references.}
\os{\newpage}
\end{description}

\exercise{Emitter follower regulator}
A Zener diode \SI{5.6}{\volt} with $R_B = \SI{1}{\kilo\ohm}$ drives the base of an npn transistor ($\beta=100$)
from $V_\mathrm{in}=\SI{12}{\volt}$. The load is $R_L=\SI{47}{\ohm}$.
\begin{talist}
\ta Estimate $\Vout$ and the load current with $V_{BE}=\SI{0.7}{\volt}$ (high current).
\ta Calculate base current and Zener current.
\ta Calculate the power of the transistor. Is a TO-92 transistor (\SI{0.6}{\watt}) sufficient?
\ta Which series resistor would a Zener shunt regulator need for the same load? Compare the idle current.
\ta Compare with Falstad for $R_L = \SI{47}{\ohm}$ and \SI{470}{\ohm}.
\end{talist}

\begin{solution}
\begin{talist}
\ta $\Vout \approx \SI{5.6}{\volt}-\SI{0.7}{\volt} = \SI{4.9}{\volt}$, $I_L = \SI{4.9}{\volt}/\SI{47}{\ohm} = \SI{104}{\milli\ampere}$.
\ta $I_B = \SI{104}{\milli\ampere}/101 = \SI{1.03}{\milli\ampere}$; $I_{R_B} = \SI{6.4}{\volt}/\SI{1}{\kilo\ohm} = \SI{6.4}{\milli\ampere}$, $I_Z = \SI{5.4}{\milli\ampere}$.
\ta $P_Q = (\SI{12}{\volt}-\SI{4.9}{\volt})\cdot\SI{104}{\milli\ampere} = \SI{0.74}{\watt}$: too much for TO-92;
    use a BD139 (TO-126) with a small heat sink.
\ta $R \le \SI{6.4}{\volt}/(\SI{104}{\milli\ampere}+\SI{5}{\milli\ampere}) = \SI{59}{\ohm}$ $\Rightarrow$ \SI{56}{\ohm}.
    Without load the Zener takes \SI{114}{\milli\ampere} (\SI{0.64}{\watt}); the series regulator only draws \SI{6.4}{\milli\ampere}.
\ta \falstad{series_regulator}: $R_L=\SI{47}{\ohm}$: $\Vout=\SI{4.89}{\volt}$, $I_L = \SI{104}{\milli\ampere}$, $V_Z=\SI{5.60}{\volt}$;
    $R_L=\SI{470}{\ohm}$: $\Vout=\SI{4.95}{\volt}$, $V_Z=\SI{5.61}{\volt}$.
\end{talist}
\end{solution}

\exercise{Adjustable regulator LM317}
\begin{talist}
\ta Design an LM317 regulator for \SI{5}{\volt} with $R_1 = \SI{240}{\ohm}$. Choose $R_2$ from E24 and calculate the actual output voltage (with $I_\mathrm{ADJ} = \SI{50}{\micro\ampere}$).
\ta Which input voltage is needed at least (dropout \SI{2}{\volt})? What is the power of the LM317 at \SI{12}{\volt} input and \SI{300}{\milli\ampere}?
\ta The LM317 is used as a current source for an LED string with \SI{20}{\milli\ampere}. Calculate the resistor.
\end{talist}

\begin{solution}
\begin{talist}
\ta $R_2 = R_1\,(5/1.25 - 1) = 3\,R_1 = \SI{720}{\ohm}$ $\Rightarrow$ E24: \SI{750}{\ohm}.
    $\Vout = \SI{1.25}{\volt}\,(1+750/240) + \SI{50}{\micro\ampere}\cdot\SI{750}{\ohm} = \SI{5.16}{\volt}+\SI{0.04}{\volt} = \SI{5.20}{\volt}$.
    (With E96: $R_2 = \SI{715}{\ohm}$ gives \SI{4.97}{\volt}.)
\ta $V_\mathrm{in} \ge \SI{7.2}{\volt}$. $P = (\SI{12}{\volt}-\SI{5.2}{\volt})\cdot\SI{0.3}{\ampere} = \SI{2.0}{\watt}$: TO-220 with heat sink.
\ta $R = \SI{1.25}{\volt}/\SI{20}{\milli\ampere} = \SI{62.5}{\ohm}$ $\Rightarrow$ E24: \SI{62}{\ohm} (\SI{20.2}{\milli\ampere}); the LM317 then needs $V_\mathrm{in} \ge V_\mathrm{LEDs} + \SI{1.25}{\volt} + \SI{2}{\volt}$.
\end{talist}
\end{solution}

%%%%%%%%%%%%%%%%%%%%%%%%%%%%%%%%%%%%%%%%%%%%%%%%%%%%%%%%%%%%%%%%%%%%%%%%%%%
\section{Transistor Current Sources}
%%%%%%%%%%%%%%%%%%%%%%%%%%%%%%%%%%%%%%%%%%%%%%%%%%%%%%%%%%%%%%%%%%%%%%%%%%%
\nsl{A bipolar transistor in its active region is a current source by
nature: its collector current depends on the base--emitter voltage, but
hardly on the collector voltage. An emitter resistor turns this into a
current source with a well-defined current. Inside integrated circuits the
current mirror is the most important building block; there is hardly an
op-amp or a comparator without several of them.\newpage}

\begin{description}
\item[{\rot\bf BJT current source}]~\\[-\lblskip]
\begin{center}
\begin{circuitikz}[scale=1, transform shape]
  \draw (0,4.5) node[vcc]{$+\SI{12}{\volt}$} -- (0,4) to[R, l_=$R_1$] (0,2) to[R, l_=$R_2$] (0,0);
  \draw (3,2) node[npn](Q){};
  \draw (0,2) -- (Q.B);
  \draw (0,2) node[circ]{} node[above right]{$V_B$};
  \draw (3,4.5) node[vcc]{$+\SI{12}{\volt}$} -- (3,4.2) to[R, l=$R_L$, i>^=$I_C$] (Q.C);
  \draw (Q.E) to[R, l=$R_E$] (3,0);
  \draw (0,0) -- (3,0);
  \draw (1.5,0) node[ground]{};
  \draw (Q.E) node[circ]{} node[right]{$V_E$};
\end{circuitikz}
\end{center}
\keyeq{BJT current source}{$\displaystyle I_C \approx I_E = \frac{V_B - V_{BE}}{R_E}\approx\frac{V_B - \SI{0.65}{\volt}}{R_E}$}
\bi
	\ite the emitter follows the base: $V_E = V_B - V_{BE}$ is fixed, so the emitter current is fixed
	\ite the load sits in the collector and may vary within the compliance range
\ei
\os{\newpage}

\item[{\rot\bf Compliance range and base voltage}]~\\[-\lblskip]
\bi
	\ite the transistor must stay active: $V_{CE} \geq$ about \SI{0.3}{\volt}, i.e.\ $V_C \geq V_E + \SI{0.3}{\volt}$
	\ite compliance: $\displaystyle R_{L,\max} = \frac{V_{CC} - V_E - \SI{0.3}{\volt}}{I_C}$; a larger load saturates the transistor and the current drops
	\ite stiff base voltage: divider current $\geq 10\,I_B$; base current error $I_B\cdot(R_1\parallelto R_2)$
	\ite temperature: $V_{BE}$ drops by \SI{2}{\milli\volt\per\kelvin}; a large $V_E$ (\SI{1}{\volt} and more) makes the current stable
\ei
\begin{center}
\begin{circuitikz}[scale=0.85, transform shape]
  \draw (0,2) node[vcc]{$V_{CC}$} -- (0,1.7) to[R, l_=$R_1$] (0,0.2) -- (0,0) to[zDo, l_=$V_Z$, *-] (0,-1.5) node[ground]{};
  \node[font=\small, align=left, anchor=west] at (0.5,0.3) {Zener};
  \draw (4,2) node[vcc]{$V_{CC}$} -- (4,1.7) to[R, l_=$R_1$] (4,0.2) -- (4,0) to[leD, *-] (4,-1.5) to[leD] (4,-3);
  \draw (4,-3) node[ground]{};
  \node[font=\small, align=left, anchor=west] at (4.5,0.3) {2 LEDs};
\end{circuitikz}
\end{center}
\nsl{The base voltage can come from a divider, from a Zener diode or from
forward-biased diodes or LEDs. The diode solutions make $V_B$ independent
of the supply voltage. Two silicon diodes ($V_B \approx \SI{1.3}{\volt}$)
have the additional advantage that their temperature coefficient cancels
that of $V_{BE}$: the voltage across $R_E$ is one diode drop and stays
constant. A red LED gives $V_B \approx \SI{1.8}{\volt}$ and
$V_E\approx\SI{1.15}{\volt}$, again with partial temperature compensation.}
\os{\newpage}

\item[{\rot\bf Example: \SI{1}{\milli\ampere} current source}]~\\[-\lblskip]
\bi
	\ite $V_{CC} = \SI{12}{\volt}$, $R_1 = \SI{10}{\kilo\ohm}$, $R_2 = \SI{2.2}{\kilo\ohm}$: $V_B = \SI{2.16}{\volt}$ (unloaded)
	\ite $R_E = \SI{1.5}{\kilo\ohm}$: $I_E = (\SI{2.16}{\volt}-\SI{0.65}{\volt})/\SI{1.5}{\kilo\ohm} = \SI{1.0}{\milli\ampere}$
	\ite loads \SI{1}{\kilo\ohm} and \SI{4.7}{\kilo\ohm}: the same current; \SI{15}{\kilo\ohm} is outside the compliance range
\ei
\falstad{bjt_current_source}
\nsl{The simulation gives \SI{1.023}{\milli\ampere} for both the
\SI{1}{\kilo\ohm} and the \SI{4.7}{\kilo\ohm} load. The difference to the
hand calculation comes from $V_{BE}$: the Falstad transistor model has
$V_{BE}=\SI{0.60}{\volt}$ at \SI{1}{\milli\ampere}, real small signal
transistors (BC547) are closer to \SI{0.65}{\volt}. With
\SI{15}{\kilo\ohm} the collector would have to go down to
$\SI{12}{\volt}-\SI{15}{\volt} < 0$; the transistor saturates and the
current falls to \SI{0.71}{\milli\ampere}. The Falstad model has no Early
effect, so in the active region the current does not depend on the load at
all. A real transistor has an output resistance of some
\SI{100}{\kilo\ohm}, increased by the emitter resistor to several
\si{\mega\ohm}.}
\os{\newpage}

\item[{\rot\bf Current mirror}]~\\[-\lblskip]
\begin{center}
\begin{circuitikz}[scale=1.1, transform shape]
  \draw (0,0) node[npn, xscale=-1](Q1){};
  \draw (3,0) node[npn](Q2){};
  \draw (Q1.B) -- (Q2.B);
  \draw (Q1.C) -- ++(0,0.4) coordinate(c1) to[R, l_=$R_\mathrm{ref}$] ++(0,2) node[vcc]{$+\SI{12}{\volt}$};
  \draw (c1) node[circ]{} -| ($(Q1.B)!0.5!(Q2.B)$) node[circ]{};
  \draw (Q2.C) -- ++(0,0.4) to[R, l=$R_L$] ++(0,2) node[vcc]{$+\SI{12}{\volt}$};
  \draw (Q1.E) -- ++(0,-0.4) coordinate(e1) -- (e1 -| Q2.E) -- (Q2.E);
  \draw ($(e1)!0.5!(e1 -| Q2.E)$) node[ground]{};
  \node[font=\small] at (-0.9,0.9) {$Q_1$};
  \node[font=\small] at (3.9,0.9) {$Q_2$};
\end{circuitikz}
\end{center}
\bi
	\ite $Q_1$ is connected as a diode: it sets $V_{BE}$ for the reference current
	\ite $Q_2$ has the same $V_{BE}$ and therefore the same collector current: $I_\mathrm{out} = I_\mathrm{ref}\,\beta/(\beta+2)$
	\ite $I_\mathrm{ref} = (V_{CC}-V_{BE})/R_\mathrm{ref}$; with \SI{11}{\kilo\ohm}: \SI{1.04}{\milli\ampere} $\to$ \SI{1.02}{\milli\ampere}
\ei
\falstad{current_mirror}
\nsl{The current mirror works only if both transistors are identical and
at the same temperature, so it is used in integrated circuits or with
matched transistor pairs (BCM847). With discrete transistors emitter
resistors of a few \SI{100}{\ohm} improve the matching. The output
current changes with the collector voltage because of the Early effect;
the Wilson mirror and cascode mirrors reduce this. Mirrors with several
output transistors distribute one reference current to many stages; by
scaling the emitter areas, the output current becomes a multiple of the
reference current.}
\os{\newpage}
\end{description}

\exercise{BJT current source for a resistive sensor}
A resistive sensor ($R_S = \SIrange{1.0}{1.4}{\kilo\ohm}$) shall be excited with \SI{1}{\milli\ampere}
by a BJT current source at $V_{CC}=\SI{12}{\volt}$ ($V_{BE}=\SI{0.65}{\volt}$, $\beta = 100$).
\begin{talist}
\ta The base voltage comes from $R_1=\SI{10}{\kilo\ohm}$, $R_2=\SI{2.2}{\kilo\ohm}$. Calculate $V_B$ with and without the base current, and choose $R_E$ from E12.
\ta Calculate the voltage across the sensor at \SI{1.0}{\kilo\ohm} and \SI{1.4}{\kilo\ohm}, and the compliance range.
\ta The base divider is replaced by two red LEDs in series ($V_F = \SI{1.8}{\volt}$ each) fed by \SI{4.7}{\kilo\ohm}. Recalculate $R_E$.
\ta Read the collector currents for \SI{1}{\kilo\ohm}, \SI{4.7}{\kilo\ohm} and \SI{15}{\kilo\ohm} in Falstad.
\end{talist}

\begin{solution}
\begin{talist}
\ta Unloaded: $V_B = \SI{12}{\volt}\cdot 2.2/12.2 = \SI{2.16}{\volt}$.
    $I_B \approx \SI{10}{\micro\ampere}$ through $R_1\parallelto R_2=\SI{1.8}{\kilo\ohm}$ lowers it by \SI{18}{\milli\volt}: $V_B=\SI{2.15}{\volt}$.
    $R_E = (\SI{2.15}{\volt}-\SI{0.65}{\volt})/\SI{1}{\milli\ampere} = \SI{1.5}{\kilo\ohm}$ (E12).
\ta $V_S = \SI{1.0}{\volt}$ and \SI{1.4}{\volt}. With $V_E = \SI{1.5}{\volt}$ and $V_{CE}\ge\SI{0.3}{\volt}$:
    $R_{L,\max} = (12-1.5-0.3)\,\si{\volt}/\SI{1}{\milli\ampere} = \SI{10.2}{\kilo\ohm}$; the sensor is far inside.
\ta $V_B = \SI{3.6}{\volt}$, $R_E = (\SI{3.6}{\volt}-\SI{0.65}{\volt})/\SI{1}{\milli\ampere} = \SI{2.95}{\kilo\ohm}$
    $\Rightarrow$ E24 \SI{3.0}{\kilo\ohm} (\SI{0.98}{\milli\ampere}). $V_B$ no longer depends on $V_{CC}$.
\ta \falstad{bjt_current_source}: $V_B = \SI{2.15}{\volt}$, $V_E = \SI{1.55}{\volt}$,
    $I_C = \SI{1.023}{\milli\ampere}$ for \SI{1}{\kilo\ohm} and \SI{4.7}{\kilo\ohm} ($V_C = \SI{7.19}{\volt}$),
    \SI{0.71}{\milli\ampere} for \SI{15}{\kilo\ohm} (saturated).
\end{talist}
\end{solution}

%%%%%%%%%%%%%%%%%%%%%%%%%%%%%%%%%%%%%%%%%%%%%%%%%%%%%%%%%%%%%%%%%%%%%%%%%%%
\section{Op-Amp Current Sources and Sensor Excitation}
%%%%%%%%%%%%%%%%%%%%%%%%%%%%%%%%%%%%%%%%%%%%%%%%%%%%%%%%%%%%%%%%%%%%%%%%%%%
\nsl{The BJT current source depends on $V_{BE}$ and $\beta$. An op-amp in
the loop removes the $V_{BE}$ error completely: it forces the voltage
across the sense resistor to equal a set voltage. Such voltage-to-current
converters are used for sensor excitation, for \SIrange{4}{20}{\milli\ampere}
current loops and as programmable current sinks. Op-amps are treated in
detail in a later chapter; here we only need the rule that an op-amp with
negative feedback makes the voltage between its inputs zero.\newpage}

\begin{description}
\item[{\rot\bf Op-amp and transistor}]~\\[-\lblskip]
\begin{center}
\begin{circuitikz}[scale=0.85, transform shape]
  \draw (0,2) node[op amp, noinv input up](A){};
  \draw (A.+) -- ++(-0.8,0) node[left]{$V_\mathrm{set}$} node[ocirc]{};
  \draw (A.out) to[R, l=\SI{1}{\kilo\ohm}] ++(1.8,0) coordinate(b);
  \draw (b) -- ++(0.3,0) node[npn, anchor=B](Q){};
  \draw (Q.C) to[R, l=$R_\mathrm{sensor}$, i<^=$I$] ++(0,2) node[vcc]{$+\SI{12}{\volt}$};
  \draw (Q.E) -- ++(0,-0.4) coordinate(e) to[R, l=$R_S$] ++(0,-1.6) node[ground]{};
  \draw (e) node[circ]{} -| ($(A.-)+(-0.5,0)$) -- (A.-);
\end{circuitikz}
\end{center}
\keyeq{Op-amp current source}{$\displaystyle V_{R_S} = V_\mathrm{set}\quad\Rightarrow\quad I_E = \frac{V_\mathrm{set}}{R_S}\qquad I = I_C = I_E\,\frac{\beta}{\beta+1}$}
\bi
	\ite $V_{BE}$ and its drift are inside the loop and drop out; the base current is lost (\SI{1}{\percent} at $\beta = 100$, a MOSFET avoids it)
\ei
\os{\newpage}

\item[{\rot\bf Floating load in the feedback path}]~\\[-\lblskip]
\begin{center}
\begin{circuitikz}[scale=1, transform shape]
  \draw (0,0) node[op amp, noinv input up](A){};
  \draw (A.+) -- ++(-1.2,0) node[ocirc]{} node[left]{$V_\mathrm{set}$};
  \draw (A.out) -- ++(0.8,0) coordinate(o);
  \draw (A.-) -| (-2,-2) coordinate(m);
  \draw (m) to[R, l_=$R_\mathrm{sensor}$, i<^=$I$] (m -| o) -- (o);
  \draw (m) node[circ]{} to[R, l_=$R_S$] ++(0,-1.8) node[ground]{};
  \draw (o) node[circ]{} node[right]{$\Vout = V_\mathrm{set}\,(1+R_\mathrm{sensor}/R_S)$};
\end{circuitikz}
\end{center}
\bi
	\ite non-inverting amplifier: the current $V_\mathrm{set}/R_S$ through $R_S$ must flow through the sensor
	\ite exact current (no base current); but neither sensor terminal is at ground
	\ite compliance: $\Vout$ must stay inside the op-amp output range
\ei
\falstad{opamp_current_source}
\os{\newpage}

\item[{\rot\bf Howland current pump (grounded load)}]~\\[-\lblskip]
\begin{center}
\begin{circuitikz}[scale=1, transform shape]
  \draw (0,0) node[op amp](A){};
  \draw (A.out) -- ++(0.6,0) coordinate(o);
  \draw (A.-) -- ++(-0.6,0) coordinate(m) -- ++(0,1) coordinate(t);
  \draw (t) to[R, l_=$R$] ++(-2.4,0) node[ground]{};
  \draw (t) node[circ]{} to[R, l=$R$] (t -| o) -- (o);
  \draw (A.+) -- ++(-0.6,0) coordinate(p) -- ++(0,-1) coordinate(q);
  \draw (q) to[R, l=$R$] ++(-2.4,0) node[ocirc]{} node[left]{$\Vin$};
  \draw (q) node[circ]{} to[R, l=$R$] (q -| o) -- (o);
  \draw (o) node[circ]{} node[right]{$\Vout$};
  \draw (q) -- ++(0,-0.4) to[R, l=$R_L$, i>_=$I_L$] ++(0,-1.8) node[ground]{};
\end{circuitikz}
\end{center}
\bi
	\ite four equal resistors: $I_L = \Vin/R$, independent of $R_L$ (within the compliance range)
	\ite the resistors must be matched well: a mismatch of \SI{0.1}{\percent} limits the output resistance
\ei
\nsl{With four equal resistors the inverting input is at $\Vout/2$.
Because the op-amp makes both inputs equal, the load voltage is
$V_L = \Vout/2$. KCL at the non-inverting input gives
$I_L = (\Vin - V_L)/R + (\Vout - V_L)/R = \Vin/R$: the load voltage
cancels. The Howland circuit is elegant but sensitive to resistor
tolerances; integrated difference amplifiers with laser-trimmed resistors
(e.g.\ INA105) are a good basis for it.}
\os{\newpage}

\item[{\rot\bf Sensor excitation: why \SI{1}{\milli\ampere}?}]~\\[-\lblskip]
\bi
	\ite resistive temperature sensor (Pt1000 like): $R(T) = R_0\,(1+\alpha T)$, $R_0 = \SI{1}{\kilo\ohm}$, $\alpha = \SI{3.85e-3}{\per\kelvin}$
	\itee \SI{0}{\celsius}: \SI{1000}{\ohm}, \SI{100}{\celsius}: \SI{1385}{\ohm}
	\ite with a constant current $I$: $V_\mathrm{sensor} = I\cdot R(T)$ is linear in $T$: \SI{3.85}{\milli\volt\per\kelvin} at \SI{1}{\milli\ampere}
	\ite more current gives more signal, but \textbf{self-heating}: $\Delta T = P/G_\mathrm{th}$
	\itee $P = I^2 R = (\SI{1}{\milli\ampere})^2\cdot\SI{1.385}{\kilo\ohm} = \SI{1.39}{\milli\watt}$; with a typical \SI{0.4}{\kelvin\per\milli\watt}: \SI{0.55}{\kelvin} error
	\ite the signal offset ($\SI{1}{\volt}$ at \SI{0}{\celsius}) is removed by a \textbf{bridge}
\ei
\nsl{The power in the sensor heats it above the temperature it should
measure. The self-heating coefficient depends strongly on the mounting:
a thin-film sensor in still air heats up by some tenths of a kelvin per
milliwatt, in flowing water ten times less. Since the power grows with the
square of the current, halving the current reduces the error to one
quarter, but also halves the signal. \SI{1}{\milli\ampere} for a
\SI{1}{\kilo\ohm} sensor (and \SI{0.3}{\milli\ampere} for a Pt1000 in
precise applications, \SI{1}{\milli\ampere} for a Pt100) are common
compromises.\newline
A current source makes the sensor voltage proportional to its resistance,
but the large constant part ($I R_0$) remains. The capstone task therefore
uses a Wheatstone bridge: a second branch with fixed resistors produces
the voltage that the sensor has at the lower end of the range, and the
instrumentation amplifier only amplifies the difference. The bridge is fed
from a stable voltage, and its resistors are chosen so that the sensor
current stays below \SI{1}{\milli\ampere}. The chapter on bridge circuits
treats this in detail.}
\os{\newpage}

\item[{\rot\bf Example: \SI{1}{\milli\ampere} for a sensor}]~\\[-\lblskip]
\bi
	\ite $V_\mathrm{set} = \SI{1.00}{\volt}$ (from a reference), $R_S = \SI{1.00}{\kilo\ohm}$ (\SI{0.1}{\percent}): $I = \SI{1.000}{\milli\ampere}$
	\ite op-amp and transistor: sensor current \SI{0.990}{\milli\ampere} (base current missing)
	\ite floating sensor in the feedback path: \SI{1.000}{\milli\ampere}, $\Vout = \SI{1}{\volt}\,(1+1.385) = \SI{2.385}{\volt}$ at \SI{100}{\celsius}
	\ite accuracy of the current = accuracy of $V_\mathrm{set}$ and $R_S$
\ei
\falstad{opamp_current_source}
\nsl{Falstad shows exactly these numbers: \texttt{VS} $=\SI{1.000}{\volt}$,
the sensor in the collector carries \SI{0.990}{\milli\ampere}, the floating
sensor \SI{1.000}{\milli\ampere} with \texttt{VOUT2} $=\SI{2.385}{\volt}$.
The \SI{1}{\kilo\ohm} resistor between op-amp and base limits the output
current of the op-amp and decouples the capacitive load of the base. For
precise sensor measurements the ratiometric principle is often even
better than a precise current: the same reference voltage that sets the
current is also used as the reference of the ADC, so its errors cancel.}
\os{\newpage}
\end{description}

%%% Local Variables:
%%% mode: latex
%%% TeX-master: "docu"
%%% End:
